%================================
% note-setup.tex
% fenglielie@qq.com 2025-09-12
%================================

\usepackage{amsmath,amsthm,amsfonts,amssymb}
\usepackage{mathtools}
\usepackage{mathrsfs}
\usepackage{bm}
\usepackage{extarrows}
\usepackage[a4paper, margin=1in]{geometry}
\usepackage{float}
\usepackage{indentfirst}
\usepackage{anyfontsize}
\usepackage{booktabs,multirow,multicol}
\usepackage[shortlabels,inline]{enumitem}
\usepackage{appendix}

\usepackage[dvipsnames]{xcolor}
\usepackage{circuitikz}
%\usepackage{graphicx}
%\graphicspath{
%    {./figure/}{./figures/}{./image/}{./images/}{./graphic/}{./graphics/}{./picture/}{./pictures/}
%}
\usepackage{subcaption}

\usepackage[ruled,linesnumbered,noline]{algorithm2e}
\usepackage{listings}
\lstdefinestyle{simpleStyle}{
    basicstyle=\ttfamily\small,
    breaklines=true,
    keywordstyle=\color{blue},
    identifierstyle=\color{black},
    stringstyle=\color{violet},
    commentstyle=\color[RGB]{34,139,34},
    showstringspaces=false,
    numbers=left,
    numbersep=2em,
    numberstyle=\footnotesize,
    frame=single,
    framesep=1em,
}
\lstset{style=simpleStyle}

\usepackage{hyperref}
\hypersetup{
    colorlinks=true,linkcolor=,urlcolor=cyan
}
\usepackage[nameinlink]{cleveref}

\renewcommand*{\proofname}{\normalfont\bfseries Proof}

\usepackage{thmtools}

%% define environments

\declaretheorem[style=plain, name=Theorem, numbered=yes, numberwithin=section, postheadhook={\phantomsection}]{theorem}
\declaretheorem[style=plain, name=Theorem, numbered=no]{theorem*}

\declaretheorem[style=plain, name=Proposition, numbered=yes, sibling=theorem, postheadhook={\phantomsection}]{proposition}
\declaretheorem[style=plain, name=Proposition, numbered=no]{proposition*}

\declaretheorem[style=plain, name=Corollary, numbered=yes, sibling=theorem, postheadhook={\phantomsection}]{corollary}
\declaretheorem[style=plain, name=Corollary, numbered=no]{corollary*}

\declaretheorem[style=plain, name=Lemma, numbered=yes, sibling=theorem, postheadhook={\phantomsection}]{lemma}
\declaretheorem[style=plain, name=Lemma, numbered=no]{lemma*}

\declaretheorem[style=plain, name=Claim, numbered=yes, sibling=theorem, postheadhook={\phantomsection}]{claim}
\declaretheorem[style=plain, name=Claim, numbered=no]{claim*}

\declaretheorem[style=definition, name=Definition, numbered=yes, numberwithin=section, postheadhook={\phantomsection}]{definition}
\declaretheorem[style=definition, name=Definition, numbered=no]{definition*}

\declaretheorem[style=definition, name=Example, numbered=yes, numberwithin=section, postheadhook={\phantomsection}]{example}
\declaretheorem[style=definition, name=Example, numbered=no]{example*}

\declaretheorem[style=definition, name=Problem, numbered=yes, numberwithin=section, postheadhook={\phantomsection}]{problem}
\declaretheorem[style=definition, name=Problem, numbered=no]{problem*}

\declaretheorem[style=remark, name=Remark, numbered=yes, numberwithin=section]{remark}
\declaretheorem[style=remark, name=Remark, numbered=no]{remark*}

\declaretheoremstyle[headfont=\color{orange!80}\bfseries, bodyfont=\normalfont, spaceabove=3pt, spacebelow=3pt]{notestyle}

\declaretheorem[style=notestyle, name=Note, numbered=yes, numberwithin=section]{note}
\declaretheorem[style=notestyle, name=Note, numbered=no]{note*}

\declaretheoremstyle[headfont=\bfseries, bodyfont=\normalfont, spaceabove=3pt, spacebelow=3pt, qed=\ensuremath{\square}]{solutionstyle}

\declaretheorem[style=solutionstyle, name=Solution, numbered=yes, numberwithin=section]{solution}
\declaretheorem[style=solutionstyle, name=Solution, numbered=no]{solution*}

% Hyperref destinations must remain unique when theorem counters reset in
% every section.  Include the section destination in each theorem anchor.
\renewcommand*{\theHtheorem}{\theHsection.\arabic{theorem}}
\renewcommand*{\theHproposition}{\theHsection.\arabic{proposition}}
\renewcommand*{\theHcorollary}{\theHsection.\arabic{corollary}}
\renewcommand*{\theHlemma}{\theHsection.\arabic{lemma}}
\renewcommand*{\theHclaim}{\theHsection.\arabic{claim}}
\renewcommand*{\theHdefinition}{\theHsection.\arabic{definition}}
\renewcommand*{\theHexample}{\theHsection.\arabic{example}}
\renewcommand*{\theHproblem}{\theHsection.\arabic{problem}}
\renewcommand*{\theHremark}{\theHsection.\arabic{remark}}
\renewcommand*{\theHnote}{\theHsection.\arabic{note}}
\renewcommand*{\theHsolution}{\theHsection.\arabic{solution}}

\usepackage[most]{tcolorbox}

\newcommand{\newtcbenvironment}[2]{
    \tcolorboxenvironment{#1}{#2, enhanced, breakable, sharp corners, boxrule=1pt}
    \tcolorboxenvironment{#1*}{#2, enhanced, breakable, rounded corners, boxrule=1pt}
}

%% define styles

\newtcbenvironment{theorem}{colframe=RoyalPurple!8, colback=RoyalPurple!8}
\newtcbenvironment{proposition}{colframe=RoyalPurple!8, colback=RoyalPurple!8}
\newtcbenvironment{corollary}{colframe=SkyBlue!8, colback=SkyBlue!8}
\newtcbenvironment{lemma}{colframe=SkyBlue!8, colback=SkyBlue!8}
\newtcbenvironment{claim}{colframe=Goldenrod!12, colback=Goldenrod!12}

\newtcbenvironment{definition}{colframe=ForestGreen!5, colback=ForestGreen!5}
\newtcbenvironment{example}{colframe=RawSienna!5, colback=RawSienna!5}
\newtcbenvironment{problem}{colframe=WildStrawberry!5, colback=WildStrawberry!5}

%% cbox
\newtcolorbox{cbox}[1][]{%
    enhanced,
    breakable,
    sharp corners,
    leftrule=2pt, rightrule=0pt, toprule=0pt, bottomrule=0pt,
    colframe=SkyBlue,
    colback=SkyBlue!8,
    #1
}

%% cover
\usepackage{titling}
\newcommand{\extrainfo}{}
\renewcommand{\extrainfo}[1]{\renewcommand{\extrainfocontent}{#1}}
\newcommand{\extrainfocontent}{}
\newcommand{\makecover}[1]{%
    \begin{titlepage}
        \newgeometry{margin=0in}
        \parindent=0pt
        \includegraphics[width=\linewidth]{#1} % size = 1280*1024
        \vfill
        \begin{center}
            \parbox{0.618\textwidth}{%
                \raggedleft{\bfseries \Huge \thetitle} \\[0.6pt]
                \rule{0.618\textwidth}{4pt} \\
            }
        \end{center}
        \vfill
        \begin{center}
            \parbox{0.618\textwidth}{%
                \raggedleft\Large
                \begin{tabular}{r}
                    \theauthor \\
                    \thedate   \\
                \end{tabular}%
            }
        \end{center}
        \vfill
        \begin{center}
            \parbox[t]{0.7\textwidth}{\centering \itshape \extrainfocontent}
        \end{center}
        \vfill
    \end{titlepage}
    \restoregeometry
    \thispagestyle{empty}
}
% USAGE
% \extrainfo{xxx}
% \makecover{/path/to/cover.png}

%hyperref
\newcommand{\refsec}[1]{\hyperref[sec:#1]{\ref*{sec:#1}节}}
\newcommand{\refsub}[1]{\hyperref[sub:#1]{\ref*{sub:#1}小节}}
\newcommand{\refeqn}[1]{\hyperref[eq:#1]{式\ref*{eq:#1}}}
\newcommand{\reffig}[1]{\hyperref[fig:#1]{图\ref*{fig:#1}}}

% 定理类引用：整体为超链接
\newcommand{\refthm}[1]{\hyperref[thm:#1]{定理\ref*{thm:#1}}}
\newcommand{\refprop}[1]{\hyperref[prop:#1]{命题\ref*{prop:#1}}}
\newcommand{\refcor}[1]{\hyperref[cor:#1]{推论\ref*{cor:#1}}}
\newcommand{\reflem}[1]{\hyperref[lem:#1]{引理\ref*{lem:#1}}}
\newcommand{\refclaim}[1]{\hyperref[claim:#1]{断言\ref*{claim:#1}}}
\newcommand{\refdef}[1]{\hyperref[def:#1]{定义\ref*{def:#1}}}
\newcommand{\refexmp}[1]{\hyperref[exmp:#1]{例\ref*{exmp:#1}}}
\newcommand{\refprob}[1]{\hyperref[prob:#1]{问题\ref*{prob:#1}}}

\crefname{theorem}{定理}{定理}
\crefname{proposition}{命题}{命题}
\crefname{corollary}{推论}{推论}
\crefname{lemma}{引理}{引理}
\crefname{claim}{断言}{断言}
\crefname{definition}{定义}{定义}
\crefname{example}{例}{例}
\crefname{problem}{问题}{问题}
